% This section will need to be drastically reworked and it is best that this is done much later in the semester when the project has come together in a different way.
% \section{Aims}
% AGN-host contamination is widely recognized as an issue that needs to be overcome to better understand the relationship and feedback mechanisms of an AGN-host system. Several methods have been developed to mitigate this contamination so that the intrinsic properties of a galaxy can be observed. Despite extensive research on mitigating AGN-related biases in SFR and stellar mass estimations, their influence on galaxy colours remains poorly understood. The prevailing approach in colour-colour analyses is to either disregard AGN entirely or exclude them without quantifying the extent of their impact. Upon reviewing published studies, we find a research gap where AGN contamination is not well quantified, and no existing literature aims to quantify the effect of AGN contamination in colour space. This underscores the need for a systematic framework to assess AGN contamination and inform decisions on their inclusion or exclusion in further studies. Broadly, this research aims to explore the effect of AGN on their host galaxy by investigating how the light from an accreting supermassive black will affect the colour space of a galaxy. Ultimately, we aim to answer the following questions:
% \begin{itemize}
% \item Does introducing an AGN component into a host galaxy influence the colourspace of galaxies?
% \item Will introducing AGN into the colourspace of galaxies cause a predictable change in a galaxy's colours?
% \item Is it possible to create a method of quantifying the change in a galaxy's colours when AGN is present?
% \end{itemize}
% \section{Hypothesis}
% Due to variations in luminosity and viewing angle, we hypothesise that AGN will significantly and measurably affect the colours of their host galaxies. This influence will manifest as contamination within color-color diagrams, potentially leading to misclassification of galaxies or an incorrect determination of galactic properties. Such issues could subsequently impact studies examining the relationship between AGN and their host galaxies and prevent further understanding of the co-evolution of these objects.
\section{Research Problem}
AGN reside in the centre of most massive galaxies, and disentangling their respective contributions to the observed light poses a significant challenge. This `AGN-host contamination' can obscure the intrinsic properties of the host galaxy, hindering our understanding of the complex interplay and feedback mechanisms within these systems. While considerable effort has been devoted to mitigating AGN-induced biases in estimating star formation rates (SFRs) and stellar masses, their impact on galaxy photometry remains poorly understood. Current analyses using colour diagnostics often neglect or exclude AGN without a quantitative assessment of their influence. This lack of understanding regarding the extent of AGN contamination in colour space presents a significant research gap.
\section{Research Aims}
This research aims to address this gap by systematically exploring the effects of AGN on their host galaxies, specifically investigating how AGN alters the observed colours of a galaxy. In doing so, we seek to answer the following questions:
\begin{itemize}
\item Does the inclusion of AGN light in galaxy colours lead to predictable changes in the associated colour diagnostic?
\item Can a method be developed to quantify the changes in a galaxy's colours due to the presence of an AGN?
\end{itemize}
By answering these questions, we aim to enhance the accuracy of galaxy characterisation by improving the robustness of colour-colour diagnostics.
\section{Hypothesis}
We hypothesise that AGN will have a significant and measurable effect on the observed colours of their host galaxies due to variations in AGN luminosity and viewing angle. This influence will manifest as contamination within colour-colour diagrams, potentially leading to:
\begin{itemize}
\item Misclassification of galaxies.
\item Inaccurate determination of galactic properties.
\end{itemize}
These issues may, in turn, impact studies investigating the relationship between AGN and their host galaxies, ultimately hindering our understanding of their co-evolution.